%oneside para una impresi\'on simple.
%titlepage define que vamos a tener una portada.
%openany para que los cap\'{\i}tulos empiecen en cualquier p\'agina y no s\'olo las impares.
%final es para compilar a modo normal. Cambiar a draft para que se compile en modo borrador. No se pone final dado que viene por defecto.
\documentclass[a4paper, 12pt, oneside, titlepage, openany]{book}

% Para el seccionado, debe colocarse al principio este package.
\usepackage{titlesec}

% Encoding + cmap (para obtener un mapeo UTF8 adecuado)
\usepackage{cmap}
\usepackage[utf8]{inputenc}
\usepackage[T1]{fontenc} % Para poder copiar y pegar el texto desde un pdf
\usepackage{lmodern}

\usepackage[backend=biber, style=iso-authoryear, language=spanish]{biblatex}
\DeclareUnicodeCharacter{202F}{~} % Para reemplazar el caracter 202F por un espacio
\bibliography{biblio.bib}
\setcounter{tocdepth}{3}
\setcounter{secnumdepth}{3}
% Esto último es para la cantidad de niveles en el índice

\DefineBibliographyStrings{spanish}{
  online   = {en línea},
  urlfrom  = {Disponible en:},
  andothers = {\mkbibemph{et\addabbrvspace al\adddot}}
}

\DeclareFieldFormat{urldate}{\mkbibbrackets{consultado\space#1}}

% Para colores
\usepackage[usenames]{color}

% Para flotar figuras y tablas
\usepackage{framed}

% Para asegurar que las figuras y tablas queden en la secci\'on correspondiente
\usepackage[section,above,below]{placeins}

% Para microtipograf\'{\i}a
\usepackage{microtype}

% Permite tama\~no de fuentes arbitrarios
\usepackage{anyfontsize}

% Itemize
\usepackage{enumitem}
\newlist{Properties}{enumerate}{2}
\setlist[Properties]{label=Propiedad \arabic*.,itemindent=*}
\renewcommand{\labelitemii}{\textasteriskcentered}
\renewcommand{\labelitemiii}{-}

\usepackage{lastpage} % Para poder contar las p\'aginas, habilita el \pageref command.
\usepackage{titleps} % Para la tabla de contenidos
\usepackage{textcase}
\usepackage{pdfpages}

% \parbox[b] es para poner texto sobre el bottom
\newpagestyle{ruled}{
	\sethead{\includegraphics[width=3.7cm]{./images/UADE_LARGE}}
			{\parbox[b]{0.43\textwidth}{\raggedright\MakeUppercase{Título de la PFI}}} % El parámetro 0.43 puede modificarse para tener más espacio en la columan del medio, y que así el título de la tesis entre mejor.
			{
				\raisebox{1.5ex}{%
					\parbox[b]{0.26\textwidth}{\raggedleft
					\begin{tabular}{r}
						Curry, Haskell \\
						Howard, William Alvin
					\end{tabular}
					}
				}
			}\headrule{}
	\setfoot[][][P\'agina \thepage\ de~\pageref*{LastPage}]{}{}{P\'agina \thepage\ de~\pageref*{LastPage}}\footrule{}
	}
\pagestyle{ruled}
% Cambiar el color de la linea divisora:
% \renewcommand\makeheadrule{\color{cyan}\rule[-.3\baselineskip]{\linewidth}{0.4pt}}
% \renewcommand\makefootrule{\color{cyan}\rule[\baselineskip]{\linewidth}{0.4pt}}

% Para que la primera página de cada capítulo no sea tratada como especial
\makeatletter
\let\ps@plain\ps@ruled
\makeatother

% M\'argenes
\usepackage{vmargin}
%\setmarginsrb{leftmargin}{topmargin}{rightmargin}{bottommargin}%
%         {headheight}{headsep}{footheight}{footskip}           %
%Lo ideal ser\'{\i}a lo siguiente
%\setmarginsrb{2.7cm}{4cm}{2.3cm}{2.5cm}{2cm}{0.5cm}{2cm}{0.5cm}
%Pero estos valores se suman, por lo que en realidad ir\'{\i}a asi:
\setmarginsrb{2.7cm}{2cm}{2.3cm}{1.5cm}{0cm}{2cm}{0cm}{1cm}
% M\'argenes, pero otra opci\'on mas acotada
% \usepackage[top=4cm,bottom=2.5cm,left=2.7cm,right=2.3cm]{geometry}

% Seccionado. El par\'ametro left incrementa el margen, lo setteo en 0.
%\titlespacing*{<command>}{<left>}{<before-sep>}{<after-sep>}
% \titlespacing{\section}{0pt}{12pt}{6pt}
%\titlespacing{\subsection}{0pt}{*0}{*0}
%\titlespacing{\subsubsection}{0pt}{*0}{*0}

%Tama\~no de letra de cap\'{\i}tulo y secciones
%\newcommand{\chapfnt}{\fontsize{16}{19}}
%\newcommand{\secfnt}{\fontsize{14}{17}}
%\newcommand{\ssecfnt}{\fontsize{12}{14}}
%\titleformat{\chapter}[display]
%{\normalfont\chapfnt\bfseries}{\chaptertitlename\ \thechapter}{20pt}{\chapfnt}
%\titleformat{\section}
%{\normalfont\secfnt\bfseries}{\thesection}{1em}{}
%\titleformat{\subsection}
%{\normalfont\ssecfnt\bfseries}{\thesubsection}{1em}{}

% Se pide numeraci\'on de tablas en n\'umeros romanos
\renewcommand{\thetable}{\thechapter.\Roman{table}}

%Espaciado antes y despu\'es de una figura respecto del texto
\setlength{\textfloatsep}{5pt}

% Para notas
\usepackage{todonotes}
\setlength{\marginparwidth}{2cm} % El margen es muy estrecho. Esto es solo para solucionar un warning
\newcommand{\Nico}[1]{\todo[color=green!30,inline]{\textbf{Nico:} #1}}
\newcommand{\NicoNI}[1]{\todo[color=green!30]{\textbf{Nico:} #1}}

%Espaciado entre referencias en la bibliograf\'{\i}a
\usepackage{etoolbox}
\patchcmd{\thebibliography}
  {\settowidth}
  {\setlength{\itemsep}{6pt}\settowidth}
  {}{}
\apptocmd{\thebibliography}
  {\small}
  {}{}

\usepackage{tocloft}
\cftsetindents{table}{0em}{4em} % Ajusta 4em seg\'un sea necesario

\usepackage{pdfcomment}
\usepackage[spanish,es-nodecimaldot]{babel}
\usepackage{csquotes}
\addto\captionsspanish{
	\renewcommand{\contentsname}{\'Indice}
	\renewcommand{\listfigurename}{Lista de Figuras}
	\renewcommand{\listtablename}{Lista de Tablas}
	\renewcommand{\tablename}{TABLA}
}

% Para que el primer parrafo de un cap\'{\i}tulo tenga identado.
\usepackage{indentfirst}
\setlength{\parindent}{36pt} %Tama\~no de la identaci\'on.

% Para tablas complejas
\usepackage{multicol,multirow}
\usepackage{array,longtable}

\usepackage{amsmath,amsfonts,amssymb,amsthm}

\allowdisplaybreaks{} % Para que el align permita dividir entre p\'aginas si lo llega a necesitar.

\usepackage{cancel}

\newenvironment{example}[1][Ejemplo]{\begin{trivlist}
	\item[\hspace{\labelsep} {\itshape{} #1}]}{\end{trivlist}}

\newenvironment{examples}[1][Ejemplos]{\begin{trivlist}
	\item[\hspace{\labelsep} {\itshape{} #1}]}{\end{trivlist}}

\newenvironment{EstadoDelArte}
    {\small\begin{center}
    \bfseries{EstadoDelArte} \end{center}}

\newenvironment{Enfoque metodologico}
    {\small\begin{center}
    \bfseries{Enfoque metodologico} \end{center}}

\newenvironment{Cronograma}
    {\small\begin{center}
    \bfseries{Cronograma} \end{center}}

% Para gr\'aficos y figuras
\usepackage{tikz}
\usepackage{graphicx}
\usepackage{wrapfig}
\usepackage{blochsphere}

% Diagramas y figuras.
\usepackage{tikz-cd}
\usepackage{quiver}

% Tablas
\newcommand\titulo[3][\scriptsize]{\rotatebox[origin=c]{90}{\parbox[t]{#2}{\centering #1{#3}}}} % Estos son los t\'{\i}tulos que van al costado de las tablas que tienen todo su contenido sin dividirse.
% \newcommand\rulestitlehalf[1]{\omit\rlap{\parbox{0.3\linewidth}{\centering\textbf{#1}}}}

\usepackage[normalem]{ulem}

\usepackage{setspace}
\onehalfspacing

% En UADE no existe la noción de capítulo en la tesis, por más que en LaTeX sí lo tratemos como tal.
% Además nos piden que el tamaño de letra de lo que para nosotros es una sección, subsección y demás, sean del mismo tamaño.
\titleformat{\chapter}
  {\bfseries\fontsize{14pt}{14pt}\selectfont}
  {\thechapter}{1em}{}
\titlespacing*{\chapter}{0pt}{12pt}{6pt}

\titleformat{\section}
  {\bfseries\fontsize{14pt}{14pt}\selectfont}
  {\thesection}{1em}{}
\titlespacing*{\section}{0pt}{12pt}{6pt}

\titleformat{\subsection}
  {\bfseries\fontsize{14pt}{14pt}\selectfont}
  {\thesubsection}{1em}{}
\titlespacing*{\subsection}{0pt}{12pt}{6pt}

\titleformat{\subsubsection}
  {\bfseries\fontsize{14pt}{14pt}\selectfont}
  {\thesubsubsection}{1em}{}
\titlespacing*{\subsubsection}{0pt}{12pt}{6pt}

% -> Acá se empieza a modificar el tamaño de los títulos autogenerados, como los de las listas de figuras y tablas
\usepackage{tocloft}

\renewcommand{\cfttoctitlefont}{\bfseries\fontsize{14pt}{16pt}\selectfont}
\setlength{\cftbeforetoctitleskip}{12pt}
\setlength{\cftaftertoctitleskip}{6pt}

% Fuente del título (14pt, negrita)
\renewcommand{\cftloftitlefont}{\bfseries\fontsize{14pt}{16pt}\selectfont}
\renewcommand{\cftlottitlefont}{\bfseries\fontsize{14pt}{16pt}\selectfont}

% Espaciado antes y después del título
\setlength{\cftbeforeloftitleskip}{12pt}
\setlength{\cftafterloftitleskip}{6pt}
\setlength{\cftbeforelottitleskip}{12pt}
\setlength{\cftafterlottitleskip}{6pt}


\usepackage[nottoc]{tocbibind}
% IMPORTANTE: En la entrega de UADE se pide en negro, pero visualmente es mejor trabajar con color
\usepackage{hyperref}
\hypersetup{
    colorlinks,
    citecolor=red
}
% Para la entrega de UADE:
%\usepackage[hidelinks]{hyperref}
%\hypersetup{
%    colorlinks=false,
%    pdfborder={0 0 0},
%    linktoc=none
%}

% -> Acá terminamos

\usepackage{mathptmx}

\DefineBibliographyStrings{spanish}{
  online   = {En línea},
  urlfrom  = {Disponible en:},
  andothers = {\mkbibemph{et\addabbrvspace al\adddot}}
}

\begin{document}
\clearpage
\thispagestyle{empty}
\includepdf[
  pages=1,
  scale=1.01, %1.17
  offset=25mm -25mm,
  pagecommand={\thispagestyle{empty}},
  noautoscale,
  trim=0 0 0 0,
  clip
]{cover/6d6a18e7-ec20-45ae-99ea-084094abeebe.pdf}
\clearpage
	\begin{titlepage}
    \centering

    {\textbf{\fontsize{18}{20}\selectfont PROYECTO FINAL DE INGENIERÍA} \par}
    \vspace{1.5cm}

    {\textbf{\fontsize{16}{18}\selectfont SmartStocker: Plataforma web para la predicción de ventas y optimización de inventario en CABA en 2025} \par}
    \vspace{0.5cm}

    {\textbf{\fontsize{14}{16}\selectfont Cano, Nicolás Martín -- LU 1147246} \par}
    {\fontsize{14}{16}\selectfont Ingeniería en Informática \par}
    \vspace{1cm}

    {\textbf{\fontsize{14}{16}\selectfont Venegas Ramírez, Gabriela Rossana -- LU 1140594} \par}
    {\fontsize{14}{16}\selectfont Ingeniería en Informática \par}
    \vspace{1.5cm}

    {\fontsize{14}{16}\selectfont Tutor: \par}
    {\textbf{\fontsize{14}{16}\selectfont Monzón, Nicolás Alberto, 
        \\ (UADE) Universidad Argentina de la Empresa, Lima 757, Cdad. Autónoma de Buenos Aires, Argentina.
        \\ (UdelaR) Universidad de la República, Av. 18 de Julio 1968, 11200 Montevideo, Departamento de Montevideo, Uruguay.
    } \par}
    \vspace{3cm}

	{\textbf{\fontsize{14}{16}\selectfont \the\year} \par}
    % {\textbf{\fontsize{14}{16}\selectfont 2025} \par} % se puede hardcodear
	% {\textbf{\fontsize{14}{16}\selectfont \today} \par} % fecha completa si es la entrega final
    \vspace{2cm}

    \includegraphics[width=0.30\textwidth]{./images/UADE}\par \vspace{1cm}
    {\textbf{\fontsize{14}{16}\selectfont UNIVERSIDAD ARGENTINA DE LA EMPRESA} \par}
    {\fontsize{14}{16}\selectfont FACULTAD DE INGENIERÍA Y CIENCIAS EXACTAS \par}
\end{titlepage}
	\chapter*{Agradecimientos}

Completar. \sout{Ejemplo de texto tachado}
	\newpage
	\chapter*{Resumen}

%% Si se quiere enviar este resumen en una p\'agina a alguien, descomentar esto.

%\begin{center}
%\textbf{\large Universidad Argentina De la Empresa}

%\textbf{Facultad de Ingenier\'{\i}a y Ciencias Exactas}

%\textbf{Departamento de Tecnolog\'{\i}a Inform\'atica}

%\end{center}

%\begin{center}
%\large Ingenier\'{\i}a en Inform\'atica
%\end{center}

%\begin{center}
%\textbf{\large Título de la PFI}
%\end{center}

%\begin{center}
%Apellido, Nombres
%Apellido, Nombres
%\end{center}

%\noindent 
Acá va el resumen. \NicoNI{Este es un comentario de margen. Tener en cuenta que el breve margen que tenemos, anula un poco su utilidad. Son libres de usarlo de todas formas}

Con el siguiente comando se imprime el interlineado en el documento pdf. \\
Interlineado simple: 12pt.\\
Interlineado esperado: 18pt.\\
Interlineado exacto de este documento: \the\baselineskip.
	\newpage
	\chapter*{Abstract}

This Final Engineering Project proposes the development of a predictive web platform for inventory management in gastronomic establishments in the Autonomous City of Buenos Aires. The main objective is to optimize inventory stock in gastronomic SMEs by 2025 by predicting sales based on a machine learning model, reducing economic losses due to overstocking or shortages, and improving operational efficiency.

The proposal consists of a web platform aimed at gastronomic SMEs, which integrates sales data from platforms like PedidosYa and uses machine learning models to estimate future demand. Unlike existing commercial solutions in the Argentine market, the system allows the incorporation of external variables (weather, holidays) and user feedback to adjust its predictions.

	\newpage
	\tableofcontents
	\chapter{Introducción}

El sector gastronómico ha sido identificado como un área estratégica dentro de la economía de la Ciudad Autónoma de Buenos Aires, destacándose por su capacidad de transformar barrios a través de la generación masiva de empleos y una variada oferta culinaria, demostrando que esta actividad no solo es fundamental para el turismo, sino también para la vida cotidiana de los residentes del país \parencite{srsur2024gastronomico}. En línea con esta relevancia, la industria gastronómica en Argentina es una de las más dinámicas y resilientes a la situación económica. A pesar de que indicadores de la Cámara Argentina de Comercio y Servicios (CAC) \parencite{cac2024consumo} determinaron que el consumo en el país cerró el 2024 con una caída acumulada de -7,4\%, en el presente el sector gastronómico experimenta un incremento constante que apunta hacia la recuperación. Pero aún cuando tanto empresarios como operadores culinarios han utilizado estrategias financieras y alternativas creativas para tratar de superar las barreras que impone la inestabilidad económica del país, sigue siendo necesario encontrar mecanismos que eleven la eficiencia operativa y así conseguir un punto de equilibrio entre la rentabilidad y la sostenibilidad.

Por lo tanto, ante el panorama presentado los distintos establecimientos gastronómicos buscan la optimización de sus operaciones, apalancándose en soluciones tecnológicas que permitan gestionar aspectos tales como la venta y facturación de sus productos, la gestión de inventario, o análisis y reportes del funcionamiento del negocio, es decir, la tecnología pasa a tener un rol preponderante, ayudando a la toma de decisiones informadas que permitan maximizar la ganancia y reducir los costos  \parencite{agrawal2023predictive}.

En el presente contexto, los avances relacionados a la tecnología predictiva, especialmente aquellos que comprenden modelos de Machine Learning, son capaces de brindar una solución innovadora. Como lo demuestra \parencite{schmidt2022mlsales}, los restaurantes tienen la capacidad de adaptarse a estas invenciones para aprovechar dichos modelos y así predecir la demanda.

A pesar de la implementación exponencial de la informática en múltiples áreas de la industria argentina, la predicción de ventas en el rubro gastronómico sigue sin ser explotada, especialmente por fuera de sistemas propietarios.

Brindarle esta posibilidad a un negocio significa disminuir enormemente la incertidumbre alrededor del cálculo de stock requerido. Esta es una tarea crítica, puesto que comprar por encima de lo necesario, en el caso de los ítems orgánicos con una fecha de vencimiento acotada, se traduce directamente en pérdidas económicas por desperdicio. Y de forma similar, comprar por debajo continúa generando un impacto financiero negativo debido a las ventas que no pueden concretarse, pero también impacta a la reputación del negocio y la fidelización del cliente \parencite{agrawal2023predictive}.

Con fundamento en consultas efectuadas a dueños y encargados de negocios, hoy en día esta es una tarea solamente realizada a partir de la experiencia personal y basada en los datos de ventas de la última semana, no considerando aspectos externos tales como la estacionalidad o el clima, cuando se encuentra demostrado que estos son determinantes a la hora de realizar predicciones en esta industria. \parencite{tanizaki2019forecasting}.

Actualmente, se encuentran dos principales líderes de mercado en Argentina, siendo estos Maxirest y Restô, los cuales a pesar de que brindan soluciones integrales para la gestión del negocio, no atacan esta problemática en particular. Se enfocan únicamente en un análisis descriptivo de las ventas o el inventario, es decir, no abordan de forma específica la problemática de predicción de almacenamiento de stock y por ende tampoco de venta.  

Maxirest es un software integral de gestión gastronómica el cual incluye funcionalidades como reportes en tiempo real, control de inventario y caja, además de ofrecer integración con plataformas de delivery como PedidosYa o Rappi, y medios de pago como MercadoPago o Payway.  

Del mismo modo, Restô es un sistema escalable de gestión gastronómica que facilita la gestión de inventario, compras, cuentas corrientes y emisión de documentos electrónicos. Igualmente, incluye múltiples módulos administrativos y contables, junto con la integración de distintos modos de pago y delivery. 

Sin embargo, ninguno de ellos comprende capacidades de predicción de demanda mediante Machine Learning, ni estimación automática de stock basado en dichas predicciones, así como tampoco permiten el ajuste dinámico del modelo a través de la retroalimentación del usuario. En consecuencia, limitan la capacidad de los negocios para la toma de decisiones basadas en estimaciones futuras, así como tampoco contemplan una personalización avanzada según el contexto de cada local.  

Ante esta situación, se propone la creación de una plataforma web que permita a un negocio, mediante el uso de un modelo de Machine Learning, obtener predicciones de ventas, y del inventario requerido para lograr satisfacerla, es decir, ir un nivel más allá de los líderes del mercado, ofreciendo un análisis predictivo mediante el uso de tecnologías emergentes basadas en Machine Learning \parencite{posch2022bayesian}.

A fin de entrenar el modelo de Machine Learning, se utilizará la información de ventas del negocio, realizando integraciones con los sistemas de POS a fin de obtenerla en tiempo real, y a su vez, se permitirá que el usuario brinde feedback sobre las predicciones, siendo este usado para continuar ajustando el modelo a las particularidades del negocio \parencite{soto2024futuro}.

La solución está orientada a PYMES gastronómicas ubicadas en CABA, que busquen reducir la incertidumbre a la hora de tomar decisiones respecto a la gestión de su stock, utilizando para esto la información de ventas.

\section{Objetivo General}

Permitir la optimización de los niveles de stock en el rubro gastronómico de Argentina en el año 2025, mediante predicción de ventas utilizando un modelo de Machine Learning, en base a la información de ventas, a fin de reducir costos y maximizar ingresos.

\noindent\textbf{Objetivos Específicos:}

\begin{itemize}
    \item Entrevistar a dueños y/o encargados de locales gastronómicos a fin de comprender el proceso de gestión de inventario y qué factores afectan a las ventas.
    
    \item Desarrollar integraciones con sistemas de plataformas de delivery, tales como PedidosYa, a fin de obtener y unificar las ventas en tiempo real.
    
    \item Desarrollar un modelo de Machine Learning que permita pronosticar ventas y extrapolar el inventario mínimo necesario para satisfacer la demanda.
    
    \item Usar fuentes externas de datos, tales como APIs meteorológicas o de feriados, a fin de enriquecer los datos a usar para la predicción.
    
    \item Permitir que el usuario defina alertas que notifique cuando un ítem del inventario queda por debajo de un umbral.
    
    \item Permitir que el usuario brinde feedback sobre estos pronósticos, a fin de permitir que el modelo se ajuste.
    
    \item Desarrollar una plataforma web donde el usuario pueda cargar los productos que vende, indicando de qué elementos del inventario se encuentran compuestos.
\end{itemize}

\section{Alcance}

El proyecto tiene como objetivo el desarrollo de una herramienta integral para la predicción de ventas y gestión de inventario en restaurantes del área de CABA, mediante la integración con sistemas POS, un modelo de Machine Learning y una interfaz web orientada al usuario gastronómico. 

El sistema contempla: 

\begin{itemize}
    \item Integración con sistemas de POS para la recopilación unificada de datos de ventas. 
    \item Desarrollo y entrenamiento de un modelo predictivo de ventas utilizando Machine Learning. 
    \item Incorporación de variables externas y feedback de usuarios en el entrenamiento del modelo. 
    \item Página web para la visualización de ventas, predicciones e inventario. 
    \item Alertas sobre los niveles de inventario. 
\end{itemize}

En el primer release se incluyen los siguientes módulos: 

\begin{itemize}
    \item Módulo de Integración POS (Módulo 1): Desarrolla interfaces para unificar la información proveniente de los sistemas de POS más relevantes en CABA. 
    \item Módulo de Predicción de Ventas (Módulo 2): Entrena un modelo de Machine Learning con datos históricos de ventas y factores externos (clima, feriados, día, etc), considerando el feedback brindado por los usuarios a fin de ajustar aún más el modelo. 
    \item Módulo Web de Visualización (Módulo 3): Permite a los usuarios visualizar ventas históricas, predicciones, niveles de inventario requeridos y brindar feedback sobre las predicciones. 
    \item Módulo de Configuración de Productos (Módulo 4): Permite definir productos, asociar sus componentes de inventario y configurar los umbrales a utilizar por las alertas.
\end{itemize}

Sin embargo, es importante acotar que existen varios aspectos que no son incluidos dentro del alcance del presente proyecto, en vista de la modalidad y características del mismo. No comprende la integración con todos los sistemas POS disponibles, ya que se concentra únicamente en los utilizados dentro del sector gastronómico de CABA. Estos no sufren modificaciones o personalizaciones que alteren su funcionamiento.

Se limita al desarrollo de una plataforma web, por lo tanto, no incluye el diseño de una aplicación móvil o de otro tipo de software adicional. Asimismo, el modelo de Machine Learning se aplica solamente a productos gastronómicos. Por lo tanto, quedan excluidos otros tipos de negocios o productos que, aunque tengan características semejantes, no forman parte de este sector.

	\subsection{Inventario en el rubro gastronómico}

Según \parencite{chopra2019supply}, se define al inventario como el conjunto de bienes o productos que una empresa mantiene en existencia con el propósito de satisfacer la demanda futura. Puede existir en forma de materias primas, productos en proceso o productos terminados, y se mantiene debido a razones como economías de escala, incertidumbre en la demanda o suministro, y variabilidad estacional.

El concepto de \emph{supply chain} se refiere al conjunto de procesos, actores y recursos involucrados en el flujo de productos, información y servicios desde los proveedores hasta el consumidor final. \parencite{chopra2019supply} definen la cadena de suministro como \guillemotleft \emph{todos los niveles involucrados, directa o indirectamente, en satisfacer la demanda del cliente}\guillemotright.

Focalizándose en el rubro gastronómico, esto implica gestionar la compra, almacenamiento y transformación de la materia prima requerida para los distintos productos a comercializar, de forma tal que estén disponibles en el momento que son requeridos, y con los niveles de calidad esperados por el cliente.

 A su vez, existen ciertas particularidades en las materias primas que conciernen a este rubro:

 \begin{itemize}
    \item La demanda es variable, ya que factores externos, tales como el clima o feriados, provocan cambios en el patrón de consumo de los clientes.

    \item Muchas de las materias primas son perecederas, con una vida útil muy corta, lo que implica que las compras a proveedores ocurren muy frecuentemente.

    \item A su vez, los pagos a proveedores tienden a ocurrir al menos una semana luego de la entrega de mercadería, lo que implica que cualquier desvío, ya sea en cantidad o tipo de materia prima requerida se traslada directamente como una pérdida económica, debido a la naturaleza perecedera de la misma.
\end{itemize}

Este proceso de reposición de inventario, siendo uno de los que ocurre con mayor frecuencia, es a su vez uno de los de mayor costo asociado, significando para uno de los entrevistados, en mayo de 2025, alrededor de 6 millones de pesos argentinos por semana en cada uno de sus locales.

A su vez, el estimar cuanto se deberia comprar es una tarea que hoy en día, debido a los costos asociados, queda en manos del personal con mayor conocimiento del negocio. Inclusive, muchas veces recae directamente en el dueño, quien debe considerar no solo las ventas que ya ocurrieron, sino también cómo los factores externos, como los feriados, el clima y el momento del mes, podrían influir en las ventas de la próxima semana.

\subsection{Aprendizaje supervisado}

El aprendizaje supervisado es una de las principales ramas de Machine Learning y se refiere al proceso mediante el cual un algoritmo aprende a realizar predicciones a partir de un conjunto de datos (dataset) etiquetado. Según \parencite{russell2022ai}, \guillemotleft{}\emph{en el aprendizaje supervisado, el agente observa pares de entrada-salida y aprende una función que mapea desde la entrada hacia la salida. Por ejemplo, las entradas podrían ser imágenes de una cámara, cada una acompañada de una salida que indica \textquotedblleft autobús\textquotedblright{} o \textquotedblleft peatón\textquotedblright{}. Esta salida se llama etiqueta. El agente aprende una función que, dada una nueva imagen, predice la etiqueta apropiada.}\guillemotright{}.

Formalmente, el aprendizaje supervisado parte de un conjunto de entrenamiento conformado por $n$ ejemplos etiquetados:

\[(x_1, y_1), (x_2, y_2), \ldots, (x_n, y_n)\]

\noindent donde cada $x_n$ es una entrada (también conocida como vector de características) y cada $y_n$ es la salida o etiqueta correspondiente. Estos ejemplos son generados por una función desconocida $f(x)$, y el objetivo del algoritmo es encontrar una función aproximada $h(x)$, denominada \emph{hipótesis}, que generalice correctamente sobre datos no vistos.

Los problemas que son abarcados por este tipo de aprendizaje son

 \begin{itemize}
    \item Clasificación: cuando las salidas $y_n$ son categorías discretas, tales como un tipo de producto o la decisión de si realizar o no una compra.

    \item Regresión: cuando las salidas son valores numéricos continuos (por ejemplo, predicción de precios, estimación de demanda o ventas futuras).
\end{itemize}

El problema de predicción de ventas abarcado en esta tesis entra en la categoría de problemas de regresión, ya que la variable a predecir es numérica continua.

\subsection{Métricas}

A la hora de implementar algún algoritmo predictivo, es fundamental tener la capacidad de medir cuán precisas son las predicciones devueltas. Es aquí, donde las métricas toman un rol preponderante, ya que permiten evaluar de distintas formas la performance de un algoritmo.

Para el caso específico de regresión, se utilizan métricas que calculan el grado de error entre los valores reales y los valores predichos por el modelo

Basándonos en \parencite{hyndman2018forecasting, james2013isl} podemos pasar a definir las tres métricas principales para modelos de regresión.

\begin{itemize}
    \item \textbf{MAE (Mean Absolute Error)}: Mide el promedio de las diferencias absolutas entre los valores predichos por el modelo y los valores reales observados, sin considerar la dirección del error. La unidad utilizada es la misma de lo que se está prediciendo, lo que facilita su interpretación, y es menos sensible a valores atípicos en comparación con otras métricas como MSE y RMSE. Se define como:

    \[
        \mbox{MAE} = \frac{1}{n} \sum_{t=1}^{n} \left| y_t - \hat{y}_t \right|
    \]

    \item \textbf{MSE (Mean Squared Error)}: Cuantifica la diferencia promedio entre los valores reales observados y las predicciones de un modelo. Se calcula como el promedio del cuadrado de los errores, es decir, la diferencia entre los valores reales y los valores predichos. Es muy sensible a valores atípicos, y su unidad es el cuadrado de la unidad de error original, lo que hace que no sea de fácil interpretación.
    
    \[
        \mbox{MSE} = \frac{1}{n} \sum_{t=1}^{n} \left( y_t - \hat{y}_t \right)^2 
    \]

    \item \textbf{RMSE – Root Mean Squared Error:} Similar al MSE, penaliza de igual forma a los valores atípicos, pero al aplicar una raíz cuadrada al error permite que este se mantenga en la misma escala que las unidades originales, solucionando el problema de interpretación del MSE.

    \[
        \mbox{RMSE} = \sqrt{ \frac{1}{n} \sum_{i=1}^{n} (y_i - \hat{y}_i)^2 }
    \]
\end{itemize}

\subsection{CatBoost}

Antes de definir el algoritmo CatBoost, es necesario definir dos conceptos claves para comprenderlo, siendo estos el concepto de función de pérdida, y de Gradient Boosting.

Una \emph{función de pérdida} es una función matemática que cuantifica el error entre la predicción del modelo ($\hat{y}$) y el valor real observado ($y$). Su objetivo es proporcionar una medida numérica del desempeño del modelo para cada ejemplo individual o conjunto de datos, la cual puede luego utilizarse para optimizar los parámetros del modelo mediante técnicas como el descenso del gradiente.

Formalmente, dada una función de predicción $f(x)$ y un valor objetivo $y$, una función de pérdida $L(y, f(x))$ devuelve un valor escalar que representa la penalización por la discrepancia entre ambos.

Definido el concepto de función de pérdida, podemos pasar a explicar el concepto de Gradient Boosting, siendo este una técnica de ensamble supervisado que construye un modelo predictivo fuerte mediante la combinación secuencial de múltiples modelos débiles, generalmente árboles de decisión, con el objetivo de corregir los errores residuales del conjunto anterior. A cada paso, el nuevo modelo es entrenado para minimizar el error del conjunto previo, utilizando métodos de optimización basados en el gradiente de una función de pérdida.

En este caso, la función utilizada por CatBoost es RMSE, cuya fórmula matemática ha sido definida previamente

Dada la definición de Gradient Boosting, podemos pasar a explicar CatBoost (Categorical Boosting), definido en \parencite{dorogush2018catboost}, como uno de los algoritmos implementados con esta técnica. Está diseñado específicamente para manejar datasets que contienen una combinación de variables numéricas y categóricas, destacándose por su capacidad para procesar variables categóricas sin necesidad de codificación previa (como one-hot o label encoding), y de reducir el overfitting que ocurre durante el entrenamiento mediante una técnica llamada Ordered Boosting. 

A su vez, es uno de los algoritmos de Machine Learning más performantes a la hora de trabajar con datasets de las características a utilizar en esta tesis.
	\chapter{Conclusi\'on}\label{conclussions}

\section{Resumen de aportes}

Completar.

\section{Trabajo futuro}

\Nico{Este es un ejemplo de comentario}

Completar.

	\addcontentsline{toc}{chapter}{Bibliograf\'{\i}a}
	\printbibliography{}

	\appendix
	\renewcommand{\appendixname}{Anexo}
\renewcommand{\thechapter}{\Alph{chapter}}

% Comando para iniciar la sección de anexos en el índice
\newcommand{\initappendix}{%
  \addcontentsline{toc}{chapter}{Anexo}%
}

% Comando para cada anexo individual
\newcommand{\anexo}[1]{%
  \refstepcounter{chapter}%
  \chapter*{Anexo \thechapter: #1}%
  \addcontentsline{toc}{section}{Anexo \thechapter: #1}%
  \markboth{Anexo \thechapter: #1}{}%
}

\initappendix

\anexo{Cronograma de Actividades}

En líneas generales, el desarrollo del proyecto ha avanzado conforme a lo planificado. No obstante, se presentó una demora en \ldots, motivo por el cual \ldots. A continuación, se presenta el cronograma actualizado con las tareas realizadas, en curso y replanificadas, que reflejan el estado actual del proyecto.

%\begin{figure}[ht]
%    \centering
%    \includegraphics[width=0.7\textheight]{./././images/Cronograma.png}
%    \caption{Cronograma de actividades.}
%\end{figure} % En las entregas que no son la entrega final, se debe tener al menos un anexo con el cronograma (Gantt)
\anexo{Resultados de la encuesta}

Con el objetivo de comprender los hábitos de consumo en locales gastronómicos y la percepción de los clientes respecto a la disponibilidad de los productos, se realizó una encuesta a más de 150 personas. El propósito de la misma fue poder identificar cómo la falta de stock y la planificación de inventario influyen directamente en la satisfacción y fidelización de los consumidores, un factor clave en el mundo gastronómico.

En cuanto a la frecuencia de consumo, el 69.8\% de los encuestados afirmó visitar locales gastronómicos o pedir por delivery al menos una vez por semana, lo que refleja una relación constante con este tipo de establecimientos.

Respecto a la disponibilidad de productos, el 47.8\%, es decir casi la mitad de los entrevistados señaló que en el último mes le sucedió que el plato o producto deseado no se encontraba disponible. Ahora bien, la falta de stock recurrente sí reflejó tener un impacto notable en la percepción del cliente ya que el 86,8\% indicó que reduciría sus visitas si un local no mantiene la disponibilidad de los platos que ofrece. Además, mientras que el 30,2\% expresó que esta situación le genera desconfianza y afecta negativamente su experiencia, el 62,92\% afirmó que le molesta al menos un poco cuando un plato no está disponible. Si se consideran ambos grupos, puede observarse que más del 93\% de los encuestados experimenta algún grado de molestia o descontento frente a la falta de disponibilidad, lo que evidencia el peso crítico de este factor en la satisfacción del cliente.

Por otra parte, el rol de la planificación de compras resultó clave: el 98,7\% de los participantes coincidió en que una mejor gestión del inventario puede mejorar significativamente el servicio. En la misma línea, el 99,4\% declaró que estaría más dispuesto a regresar a un local que siempre mantenga su menú disponible y con calidad constante, y el 98,7\% lo recomendaría más a terceros.

Finalmente, la encuesta mostró que la falta frecuente de platos afecta directamente la reputación de los locales (71,7\% lo considera un factor muy relevante) y entre las respuestas abiertas, los clientes destacaron como principal fuente de satisfacción la combinación de calidad, disponibilidad, buena atención y relación precio-calidad, confirmando la importancia de un sistema que permita optimizar el control de insumos. % En el caso de tener en cuestas tiene que haber una transcripción de los resultados
\anexo{Transcripción de la entrevista a Alan Turing}

% En el Alan Turing Internet Scrapbook se menciona que Turing ofreció una charla en la radio de la BBC en 1951 
% titulada "Can Computers Think?", y participó en una discusión radiofónica en 1952. 
% Aunque no hay grabaciones conocidas, los guiones (scripts) originales sí están disponibles en el Turing Archive.
% De todas formas, me pareció anecdótico pedirle a un LLM cómo "piensa" que sería esta charla.
% Modelo: GPT5
% Prompt: Piensa cómo sería una entrevista a Alan Turing sobre el paper "Could a computer think?"
\begin{description}[leftmargin=0cm, labelsep=0.5cm]

    \item[\textbf{Entrevistador:}] Profesor Turing, su artículo plantea la pregunta “¿Podría pensar una computadora?”. ¿Cuál diría que es el núcleo de su argumento?

    \item[\textbf{Turing:}] No se trata tanto de si las máquinas pueden pensar, sino de qué entendemos por “pensar”. Propongo desplazar la pregunta hacia algo más verificable: ¿puede una máquina imitar exitosamente el comportamiento inteligente humano al punto de engañar a un observador? Eso es lo que llamo el juego de imitación.

    \item[\textbf{Entrevistador:}] Muchos críticos afirman que las máquinas nunca podrán “pensar” porque carecen de conciencia o emociones.

    \item[\textbf{Turing:}] La objeción de la conciencia es interesante, pero también se aplica a otros seres humanos: no puedo probar que usted sea consciente. Solo puedo observar su comportamiento. Por tanto, es más práctico juzgar la inteligencia por lo que un sistema hace, no por lo que supuestamente es.

    \item[\textbf{Entrevistador:}] En el artículo usted responde a varias objeciones, desde las religiosas hasta las matemáticas. ¿Hay alguna que considere la más difícil de refutar?

    \item[\textbf{Turing:}] Tal vez la objeción de la “infinita variedad” —la idea de que el comportamiento humano es demasiado variado para ser reproducido por reglas finitas. Sin embargo, las matemáticas nos muestran que sistemas finitos pueden generar gran complejidad. La biología también es un buen ejemplo: el código genético es finito, pero produce toda la diversidad de la vida.

    \item[\textbf{Entrevistador:}] Usted menciona que las máquinas podrían aprender. En 1950 eso sonaba muy futurista.

    \item[\textbf{Turing:}] Sí, creo que la programación explícita no es el único camino. Una máquina podría ser diseñada para modificar sus propias instrucciones y mejorar con la experiencia. Sería, en esencia, un aprendiz.

    \item[\textbf{Entrevistador:}] Si mirara hacia el futuro, ¿cómo imagina que evolucionará esta idea?

    \item[\textbf{Turing:}] Confío en que, hacia finales del siglo, las máquinas serán capaces de participar en el juego de imitación con éxito. Tal vez aún habrá debate sobre si “piensan”, pero la práctica hará irrelevante la pregunta: simplemente conviviremos con entidades que actúan inteligentemente.

\end{description}

% Esta entrevista fue real, pero acá dejo una versión acortada
\anexo{Transcripción de la entrevista a Richard Feynman}
\begin{description}[leftmargin=0cm, labelsep=0.5cm]
    \item[\textbf{Interviewer:}] If you get hold of two magnets, and you push them, you can feel this pushing between them. Turn them around the other way, and they slam together. Now, what is it, the feeling between those two magnets?

    \item[\textbf{Feynman:}] What do you mean, "What's the feeling between the two magnets?"

    \item[\textbf{Interviewer:}] There's something there, isn't there? The sensation is that there's something there when you push these two magnets together.

    \item[\textbf{Feynman:}] Listen to my question. What is the meaning when you say that there's a feeling? Of course you feel it. Now what do you want to know?

    \item[\textbf{Interviewer:}] What I want to know is what's going on between these two bits of metal?

    \item[\textbf{Feynman:}] They repel each other.

    \item[\textbf{Interviewer:}] What does that mean, or why are they doing that, or how are they doing that? I think that's a perfectly reasonable question.

    \item[\textbf{Feynman:}] Of course, it's an excellent question. But the problem, you see, when you ask why something happens, how does a person answer why something happens? For example, Aunt Minnie is in the hospital. Why? Because she went out, slipped on the ice, and broke her hip. That satisfies people. It satisfies, but it wouldn't satisfy someone who came from another planet and who knew nothing about why when you break your hip do you go to the hospital. How do you get to the hospital when the hip is broken? Well, because her husband, seeing that her hip was broken, called the hospital up and sent somebody to get her. All that is understood by people. And when you explain a why, you have to be in some framework that you allow something to be true. Otherwise, you're perpetually asking why. Why did the husband call up the hospital? Because the husband is interested in his wife's welfare. Not always, some husbands aren't interested in their wives' welfare when they're drunk, and they're angry.
    \ldots
\end{description} % En el caso de tener entrevistas, se deben poner sus transcripciones
	\newpage
	\listoffigures
	\newpage
	\listoftables
\end{document}
